\documentclass[11pt]{article}
\usepackage[utf8]{inputenc}
\usepackage[T1]{fontenc}
\usepackage{graphicx}
\usepackage{booktabs}
\usepackage{amsmath}
\usepackage{amssymb}
\usepackage[margin=1in]{geometry}
\usepackage{caption}
\usepackage{array}
\begin{document}
\title{Hearthsalt lozenge use and dusk-spot in one adult cohort}
\author{Mira Pell and Owen Silt}
\maketitle
\begin{abstract}
Records from Pell Kiln Room and Silt Ledger Room were reviewed for Hearthsalt lozenge use.
Of {{v:n_identified}} records identified, {{v:n_excluded}} were set aside and {{v:n_analyzed}} were analyzed,
{{v:n_exp}} with a lozenge and {{v:n_unexp}} without one.
Dusk-spot was recorded for {{v:ev_exp}} exposed adults ({{v:risk_exp_pct}}\%)
and {{v:ev_unexp}} unexposed adults ({{v:risk_unexp_pct}}\%).
The risk difference was {{v:rd_pp}} percentage points (95\% CI {{v:rd_lo}} to {{v:rd_hi}}).
The risk ratio was {{v:rr}} (95\% CI {{v:rr_lo}} to {{v:rr_hi}})
and the odds ratio was {{v:or_est}} (95\% CI {{v:or_lo}} to {{v:or_hi}}).
\end{abstract}
\section{Introduction}
Dusk-spot is a pale plaque that becomes visible after sunset \cite{pell2016}.
Some adults at Pell Kiln Room and Silt Ledger Room were issued a Hearthsalt lozenge during the index year, and others were not \cite{silt2018}.
A vellum mark on the forearm was written down at the first visit and was named in advance as the subgroup split \cite{harbor2019}.
The report uses a record-flow count, one outcome table, and Wald intervals for the risk difference and for the risk ratio and odds ratio on the log scale \cite{cedar2017,moss2020,quill2021}.
Lozenge use was taken from the room record. It was not assigned by the investigators.
\section{Methods}
Adults with a kiln-work record at Pell Kiln Room or Silt Ledger Room were identified.
A record was set aside when the exposure field was missing, when the dusk-spot field was missing, or when the adult was below the age cutoff used by the rooms.
Exposure was a Hearthsalt lozenge issued in the index year. The comparison group had no lozenge issued.
The outcome was a dusk-spot written in that same year.
The prespecified subgroup split the analyzed records by a vellum mark, present or absent.
The exposed risk $p_1$ is the exposed event count $a$ divided by the exposed group size $n_1$.
The unexposed risk $p_0$ is the unexposed event count $b$ divided by the unexposed group size $n_0$.
The risk difference is $p_1-p_0$.
Its Wald standard error is the square root of $p_1(1-p_1)/n_1+p_0(1-p_0)/n_0$, and the interval is that difference plus or minus 1.96 times the standard error.
The risk ratio is $p_1/p_0$.
Its standard error on the log scale is the square root of $1/a-1/n_1+1/b-1/n_0$, and the bounds are the exponential of the log risk ratio plus or minus 1.96 times that standard error.
The odds ratio is $(a/(n_1-a))/(b/(n_0-b))$.
Its standard error on the log scale is the square root of $1/a+1/(n_1-a)+1/b+1/(n_0-b)$, and the bounds are the exponential of the log odds ratio plus or minus 1.96 times that standard error.
A group risk uses the same multiplier and the Wald standard error given by the square root of $p(1-p)/n$.
Percentages and the risk difference, in percentage points, are rounded half up to one decimal place.
Risk ratios, odds ratios, and their interval bounds are rounded half up to two decimal places.
\section{Results}
Of {{v:n_identified}} records identified, {{v:n_ex_exposure}} lacked an exposure field, {{v:n_ex_outcome}} lacked a dusk-spot field, and {{v:n_ex_age}} were below the age cutoff.
Those three counts sum to {{v:n_excluded}}.
The analyzed set held {{v:n_analyzed}} adults, {{v:n_exp}} exposed and {{v:n_unexp}} unexposed.
\falsifytable{T1}
Dusk-spot occurred in {{v:ev_exp}}/{{v:n_exp}} ({{v:risk_exp_pct}}\%) exposed adults and in {{v:ev_unexp}}/{{v:n_unexp}} ({{v:risk_unexp_pct}}\%) unexposed adults.
The complementary counts were {{v:none_exp}} and {{v:none_unexp}}.
The exposed risk was {{v:risk_exp_pct}}\% (95\% CI {{v:risk_exp_lo}} to {{v:risk_exp_hi}}) and the unexposed risk was {{v:risk_unexp_pct}}\% (95\% CI {{v:risk_unexp_lo}} to {{v:risk_unexp_hi}}).
The risk difference was {{v:rd_pp}} percentage points (95\% CI {{v:rd_lo}} to {{v:rd_hi}}).
The risk ratio was {{v:rr}} (95\% CI {{v:rr_lo}} to {{v:rr_hi}}) and the odds ratio was {{v:or_est}} (95\% CI {{v:or_lo}} to {{v:or_hi}}).
\falsifytable{T2}
Where the vellum mark was present, dusk-spot occurred in {{v:ev_mark_exp}}/{{v:n_mark_exp}} ({{v:pct_mark_exp}}\%) exposed adults and in {{v:ev_mark_unexp}}/{{v:n_mark_unexp}} ({{v:pct_mark_unexp}}\%) unexposed adults.
Where the mark was absent, the counts were {{v:ev_plain_exp}}/{{v:n_plain_exp}} ({{v:pct_plain_exp}}\%) and {{v:ev_plain_unexp}}/{{v:n_plain_unexp}} ({{v:pct_plain_unexp}}\%).
The two strata add to the exposed group, the unexposed group, and the event counts.
\falsifytable{T3}
\falsifyfigure{F1}
\section{Discussion}
{{v:ph_concl_claim}}
{{v:ph_hedge}}
The groups were not formed by a chance mechanism, and kiln hours, mark status, and night-shift work were not held fixed, so unmeasured differences could explain the contrast \cite{quill2021}.
{{v:ph_limitation}}
The same four cells determine the risks, the risk difference, the risk ratio, and the odds ratio reported above \cite{moss2020,cedar2017}.
\falsifybibliography
\end{document}
